\chapter{REVISI\'ON CR\'ITICA DE LA LITERATURA}

Este capítulo revisa críticamente el estado del arte en las líneas temáticas que estructuran esta 
investigación: estudios comparativos de arquitecturas de visión bajo condiciones
federadas heterogéneas, el papel de la arquitectura como variable de diseño en FL, la integración
de modelos basados en State Space Models (SSM/Mamba) y de Vision GNNs en sistemas federados,
y los enfoques de análisis por capas que fundamentan las métricas de drift adoptadas.
El objetivo no es enumerar estudios, sino identificar tendencias, vacíos y la cadena de conocimiento
que justifica las decisiones metodológicas de esta tesis.


\section{Estudios comparativos de arquitecturas de visión en FL}

La pregunta de si la elección de arquitectura afecta de forma sistemática el comportamiento de los
modelos bajo condiciones no-IID fue formulada con rigor por primera vez por Qu et al.\
\cite{qu2022rethinking}, en lo que constituyó la primera investigación empírica exhaustiva de
múltiples familias arquitectónicas bajo distintos algoritmos federados, \textit{benchmarks} reales
y particiones heterogéneas.
Sus experimentos, que abarcaron CNNs (ResNet, EfficientNet) y Vision Transformers (ViT, Swin)
sobre CIFAR-10, Retina y CelebA, evidenciaron que las arquitecturas basadas en autoatención eran
sistemáticamente más robustas ante \textit{distributional shifts} no-IID, con mejoras de hasta 37
puntos porcentuales sobre las CNNs equivalentes bajo alta heterogeneidad.
Los autores argumentaron que el mecanismo de autoatención global, al no depender de estadísticas
de mini-batch locales para su normalización, mitigaba de forma inherente la divergencia entre
actualizaciones de distintos clientes.
Sin embargo, el trabajo no desagregaba las fuentes de esa ventaja: no se distinguía entre el
efecto del sesgo inductivo (convolución local frente a atención global) y el efecto de la estrategia
de normalización (BatchNorm en CNNs frente a LayerNorm en ViTs), ni se proporcionaba ningún
análisis por capa que permitiera localizar dónde se concentraba la divergencia.

En la misma dirección, Zuo et al.\ \cite{zuo2022empirical} realizaron una comparación temprana
entre CNNs ligeras y variantes compactas de ViT bajo restricciones de memoria y tiempo propias
de dispositivos de borde, confirmando que los ViTs entrenados desde cero con imágenes de baja
resolución podían superar a sus homólogos CNN en precisión global tanto en escenarios IID como
no-IID con un coste computacional comparable.
No obstante, el estudio operó exclusivamente a nivel de métricas de rendimiento final, sin
descomponer el comportamiento por capa ni explorar el efecto de la normalización.

Las comparaciones más recientes han ido ampliando progresivamente el espectro de arquitecturas
evaluadas.
Büyüktaş et al.\ \cite{buyuktacs2024transformer} examinaron arquitecturas del estilo MetaFormer
---MLP-Mixer, ConvMixer y PoolFormer--- frente a ResNet-50 en clasificación multietiqueta de
teledetección bajo FL con FedAvg y MOON.
Bajo alta heterogeneidad, PoolFormer superó a ResNet-50 en 12.53 puntos de F1, mientras que MOON
apenas mejoró a las arquitecturas Transformer, lo que sugería que su robustez intrínseca ante el
drift reducía el beneficio marginal de los algoritmos de regularización contrastiva.
Li et al.\ \cite{li2025comparison} obtuvieron conclusiones análogas al comparar ViT-Base con
ResNet-50 en reconocimiento de acciones humanas bajo FL: el ViT federado alcanzó 87.30\% de
precisión frente al 79.29\% del ResNet, con una brecha de generalización cuatro veces menor,
atribuida al desarrollo de capacidad de abstracción transversal a los clientes a través del
mecanismo de autoatención.

Jiménez-Gutiérrez et al.\ \cite{jimenez2503thorough} aportaron una de las evaluaciones más
sistemáticas sobre el impacto diferencial de los tipos de heterogeneidad no-IID.
Mediante la Distancia de Hellinger como métrica unificada, demostraron que el sesgo de etiquetas
producía caídas de precisión de entre 10 y 40 puntos sobre la línea base centralizada, con umbrales
críticos de degradación a $\text{HD} > 0.5$ y $\text{HD} > 0.75$, mientras que los sesgos de
cantidad y características resultaban prácticamente inocuos.
Los modelos de transferencia como EfficientNet-B0 mostraron mayor sensibilidad al sesgo de etiquetas
que las CNNs entrenadas desde cero, y ninguno de los cinco algoritmos de agregación evaluados
---FedAvg, FedProx, MOON, POC y selección aleatoria--- logró reducir la brecha de forma sustancial
bajo alta heterogeneidad.
Este resultado constituye un punto de referencia empírico directo para la presente investigación:
los valores $\alpha \in \{1000.0, 1.0, 0.3, 0.1, 0.03\}$ empleados en las particiones de Dirichlet
se corresponden aproximadamente con los niveles HD analizados en dicho trabajo, y la incapacidad
algorítmica para resolver el problema bajo alta heterogeneidad refuerza la motivación de examinar
la arquitectura como variable independiente.
Sin embargo, el estudio se limitó exclusivamente a CNNs y no incorporó análisis por capas,
familias Transformer, SSMs ni Vision GNNs.

Estudios comparativos más recientes que incluyen SSMs han aportado evidencia complementaria,
aunque en entornos centralizados.
Biswas et al.\ \cite{biswas2025comparative} compararon diecisiete modelos de tres familias
---CNN, Transformer y SSM (Vision Mamba)--- en tres dominios de imágenes, concluyendo que los
modelos CNN alcanzaban mayor precisión bruta mientras que los basados en Mamba presentaban la
menor huella de memoria y la convergencia más rápida.
Yanar et al.\ \cite{yanar2025comparative} llegaron a una conclusión más matizada en el dominio
de rayos X de tórax: los modelos híbridos CNN-Transformer (ConvFormer, CaFormer) y EfficientNet
superaron consistentemente a las variantes de Mamba (MedMamba, VMamba) en 14 categorías de
patologías, advirtiendo que la competitividad de los SSMs no es universal y depende del dominio.
Ambos estudios se desarrollaron en entornos centralizados, sin federación ni análisis de drift,
por lo que sus conclusiones no son directamente transferibles a escenarios no-IID.

En el dominio de imagenología cerebral, varios estudios han aplicado arquitecturas de visión sobre
el dataset Brain Tumor MRI.
Trabajos como los de Chen et al.\ \cite{chenfederated} y Lai et al.\ \cite{lai2025advancing}
realizaron el entrenamiento de forma centralizada.
Quienes sí incorporaron FL ---Anoosha y Seetharamulu \cite{anoosha2025federated}, Zhou et al.\
\cite{zhou2024distributedfederatedlearningbaseddeep} y Appasami y Savarimuthu
\cite{appasami2025federated}--- lo hicieron empleando exclusivamente variantes CNN, sin realizar
ningún análisis de comportamiento por capas ni comparación entre familias arquitectónicas.
Ninguno de estos estudios abordó la dinámica de drift paramétrico o representacional a lo largo
de las rondas de comunicación.

En conclusión, los trabajos previos presentan las siguientes limitaciones: primero, las evaluaciones se centran predominantemente en métricas de rendimiento
final, sin medir los procesos internos que las generan; segundo,
el espacio arquitectónico comparado raramente ha superado dos familias de modelos, ignorando sistemáticamente los SSMs
y los Vision GNNs en contextos federados; y tercero, la heterogeneidad ha sido poco estudiada como un parámetro controlado que permita comparaciones rigurosas.
La presente investigación responde a las tres limitaciones de forma directa.


\section{Arquitectura como variable de diseño en FL}

El trabajo de Pieri et al.\ \cite{pieri2023handlingdataheterogeneityarchitectural} representa el
estudio de mayor alcance en términos de diversidad arquitectónica evaluada dentro de un marco
federado.
En él se compararon 19 modelos pertenecientes a cinco familias ---CNNs, Transformers, MLP-Mixers,
modelos híbridos y MetaFormers--- sobre cuatro conjuntos de datos FL con distintos tipos de
heterogeneidad, evaluando además el efecto ortogonal de cuatro algoritmos de optimización:
FedAvg, FedAVGM, FedProx y SCAFFOLD.
Los autores reportaron que las arquitecturas tipo MetaFormer (CAFormer, ConvFormer, PoolFormer)
exhibían la mayor robustez ante distribuciones heterogéneas, y que el cambio de arquitectura
producía mejoras más consistentes y de mayor magnitud que la elección del algoritmo de agregación.
El hallazgo más relevante para la presente investigación fue la ablación de normalización: al
sustituir BatchNorm por LayerNorm en modelos CNN, la brecha de rendimiento entre CNNs y
Transformers se reducía de forma considerable, sugiriendo que BatchNorm ---y no el sesgo
inductivo convolucional \textit{per se}--- era el factor determinante de la inferioridad de las
CNNs bajo FL heterogéneo.
Sin embargo, el trabajo presentó dos vacíos fundamentales: no incluyó ninguna arquitectura SSM
ni Vision GNN, y no incorporó métricas de drift paramétrico o representacional por capa.

Xu et al.\ \cite{xu2023fedconv} ofrecieron una perspectiva complementaria: mediante ablación
sistemática sobre componentes micro-arquitectónicos de una CNN pura, demostraron que con
modificaciones dirigidas ---función de activación SiLU, entrenamiento libre de normalización con
\textit{Adaptive Gradient Clipping}, \textit{stem} convolucional solapado y \textit{kernels}
de tamaño 9--- la arquitectura resultante (FedConv) superaba a ViT-Small en precisión FL bajo
alta heterogeneidad.
Este hallazgo matizó la conclusión de Qu et al.: la inferioridad de las CNNs no era inherente a
la convolución, sino que derivaba de decisiones micro-arquitectónicas corregibles, siendo la
eliminación de BatchNorm la modificación de mayor impacto individual.
No obstante, FedConv evaluó únicamente métricas de precisión y eficiencia de comunicación, sin
medir la divergencia de pesos por capa ni la alineación de gradientes entre clientes.

Siomos et al.\ \cite{siomos2024architecture} profundizaron en la dirección
\textit{normalization-free} mediante ANFR, una arquitectura que combina \textit{Scaled Weight
Standardization} (SWS) con mecanismos de atención de canal.
Al estandarizar los pesos en lugar de las activaciones, ANFR eliminó la dependencia del paso
hacia adelante de las estadísticas de mini-batch del cliente, logrando coherencia de los
descriptores de canal entre participantes heterogéneos.
Evaluada bajo siete algoritmos de agregación y cinco conjuntos de datos, ANFR alcanzó el mejor
rendimiento en todas las configuraciones, incluyendo escenarios de privacidad diferencial donde
otras estrategias de normalización degradaban fuertemente.
El trabajo confirmó que BatchNorm era el principal vector de inestabilidad federada en CNNs,
proporcionando base empírica sólida para la ablación de normalización planteada en esta tesis.
Sin embargo, como en el caso de Pieri et al. \cite{pieri2023handlingdataheterogeneityarchitectural}, el estudio se restringió a CNNs y no realizó
mediciones de drift por capa, gradiente o representación.

En conjunto, existe evidencia que la arquitectura es un eje de diseño primario en FL,
potencialmente más determinante que la elección del algoritmo de agregación, e identificó a
BatchNorm como el principal confusor entre el efecto del sesgo inductivo y el de la normalización.
No obstante, todos los estudios midieron el impacto arquitectónico exclusivamente a través de
precisión final, ninguno trazó la evolución del drift con granularidad por capa, y ninguno
extendió el análisis a las familias SSM o Vision GNN.


\section{State Space Models y Vision Mamba en FL}

Los modelos de espacio de estados selectivos irrumpieron en el dominio de la visión por computadora
con Vision Mamba (Vim) \cite{zhu2024vision}, que introdujo un bloque SSM bidireccional capaz de
procesar secuencias de parches de imagen con complejidad lineal en memoria y tiempo.
En comparación directa con DeiT, Vim-Tiny superó al Transformer equivalente en ImageNet-1K en
3.9 puntos de \textit{top-1 accuracy}, alcanzando una velocidad 2.8$\times$ mayor y un consumo
de memoria 86.8\% inferior a resoluciones elevadas.
El diseño emplea LayerNorm al inicio de cada bloque ---en contraste con la BatchNorm de las CNNs
clásicas---, lo que lo sitúa, desde la perspectiva de la normalización, en una posición análoga a
los ViTs.
La naturaleza secuencial del SSM implica, no obstante, un mecanismo de integración de la
información fundamentalmente distinto al de la autoatención: mientras los ViTs establecen
dependencias globales simultáneas entre todos los parches, los SSMs integran la información a
través de transiciones de estado selectivas dependientes del contenido, cuyo comportamiento bajo
distribuciones de clases heterogéneas entre clientes no había sido caracterizado.

La idoneidad de los SSMs para tareas de visión fue matizada por Yu y Wang \cite{yu2025mambaout},
quienes demostraron que MambaOut ---un modelo que elimina el componente SSM reemplazándolo por
convoluciones \textit{depthwise} con \textit{kernel} grande--- superaba a todas las variantes
visuales de Mamba en clasificación de ImageNet-1K, aunque no en tareas densas.
Los autores argumentaron que la clasificación estándar no satisface las condiciones de secuencia
larga ni de procesamiento autorregresivo para las cuales el SSM fue diseñado.
En el entorno federado, donde el mini-batch local puede ser reducido y la distribución de clases
local puede ser muy sesgada, no está claro \textit{a priori} si el mecanismo SSM aportará ventajas
o desventajas en términos de drift paramétrico y representacional.

En el dominio médico, MedMamba \cite{yue2024medmamba} y la revisión de Bansal et al.\
\cite{bansal2024comprehensive} documentaron que las variantes híbridas SSM-CNN ofrecen rendimiento
competitivo en clasificación de imágenes de múltiples modalidades, aunque en entornos centralizados.
La integración de Mamba en FL fue abordada por primera vez por Wu \cite{wu2024mamba}, quien
demostró mayor precisión y eficiencia de entrenamiento respecto a \textit{baselines} convolucionales,
concluyendo que los SSMs eran viables para aplicaciones distribuidas y sensibles a la privacidad.
No obstante, el estudio no exploró particiones no-IID sistematizadas ni ofreció análisis de drift
por capa.
Trabajos posteriores como los de Chen \cite{chen2024mamaba} y Guo et al.\ \cite{guo2026mamba}
aplicaron Mamba en FL a dominios específicos ---detección de rayos X y diagnóstico fotovoltaico,
respectivamente--- con el mismo foco en precisión final y sin estudiar los mecanismos de divergencia
paramétrica inducidos por la heterogeneidad.
La comprensión del comportamiento de los SSMs ante condiciones no-IID controladas, con análisis
por capa y a lo largo de las rondas de comunicación, permanece como un vacío no abordado.


\section{Vision GNNs en FL}

Los modelos de redes neuronales en grafos se originaron como herramientas para el procesamiento
de datos estructurados de forma irregular \cite{Gori2005ANM, Scarselli:2009ku}.
Han et al.\ \cite{han2022vision} introdujeron Vision GNN (ViG), la primera arquitectura de visión
pura basada en grafos para clasificación de imágenes a escala de ImageNet, demostrando que el
mecanismo de propagación de mensajes sobre un grafo de parches podía aprender representaciones
competitivas con las obtenidas por CNNs y ViTs.
El sesgo inductivo de ViG es cualitativamente distinto: mientras las CNNs explotan la localidad
espacial mediante filtros compartidos y los ViTs establecen dependencias globales a través de la
autoatención, ViG construye dinámicamente una topología de conectividad entre regiones de la imagen
y refina las representaciones nodales iterativamente, lo que puede resultar en una sensibilidad
diferente a la distribución de clases local de cada cliente bajo federación.

La interacción entre GNNs y FL fue sistematizada por He et al.\ \cite{DBLP:journals/corr/abs-2104-07145},
quienes presentaron FedGraphNN como la primera plataforma de \textit{benchmarking} para FL con
GNNs sobre datos nativamente graficados, abarcando 36 conjuntos de datos en siete dominios.
Sus experimentos revelaron que el rendimiento federado era sistemáticamente inferior al centralizado
y que la arquitectura GNN con mejor desempeño centralizado no mantenía necesariamente su ventaja
bajo federación ---resultado que anticipa la relevancia de estudiar el efecto moderador de la
arquitectura también para modelos basados en grafos.
Sin embargo, FedGraphNN operó exclusivamente sobre datos estructurados como grafos (moléculas,
redes de citas, sistemas de recomendación), sin considerar \textit{backbones} de visión sobre
imágenes.
Plataformas más recientes como FedGraph \cite{yao2024fedgraph} ampliaron la infraestructura
disponible pero mantuvieron el foco en eficiencia y seguridad del software, sin estudiar los
mecanismos de drift paramétrico o representacional.

La revisión sistemática de Shaikh y Samet \cite{shaikh2025federated} proporcionó una taxonomía
de la heterogeneidad en FedGNNs e identificó la heterogeneidad estructural/topológica como el
tipo más descuidado en la literatura, confirmando la ausencia de métricas estandarizadas para
cuantificar la no-IID-ness en grafos.
Más relevante para esta tesis, el survey no incluyó ningún estudio que empleara el \textit{backbone}
de Vision GNN en tareas de clasificación de imágenes bajo FL.
El trabajo de Cui et al.\ \cite{cui2025fedvig} fue el primero en integrar explícitamente una
arquitectura Vision GNN en un marco federado (FedVIG), empleando perturbación de datos mediante
modelos de difusión y aprendizaje contrastivo estructural para mejorar la generalización y la
privacidad.
Sin embargo, el trabajo no estudió el comportamiento del modelo ante particiones no-IID
sistematizadas, y su foco primario fue la privacidad, no la caracterización del drift.
Estudios que emplearon GNNs sobre el dataset de tumores cerebrales, como los de Gürsoy y Kaya
\cite{gursoy2024brain} o Elhamzi y Mokni \cite{elhamzi2026braingraphnet}, lo hicieron en entornos
centralizados con modelos híbridos o variantes de GCN ---no el \textit{backbone} de Vision GNN---
y sin ninguna consideración federada.

En consecuencia, la literatura no contiene ningún estudio que analice el comportamiento de Vision
GNN bajo datos heterogéneos en un entorno federado, ni que compare Vision GNN con Vision Mamba
en ese contexto.
Este vacío ---la ausencia de análisis de drift para Vision GNN bajo FL y la inexistencia de
comparaciones entre las cuatro familias (CNN, ViT, SSM, GNN) bajo condiciones no-IID controladas---
representa la brecha central que aborda la presente investigación.


\section{Análisis por capas y métricas de drift}

La evidencia empírica que motivó el análisis granular del drift provino inicialmente de trabajos
que utilizaron el Centered Kernel Alignment (CKA) como herramienta de diagnóstico representacional.
Son et al.\ \cite{son2022comparisons} fueron los primeros en aplicar CKA en un contexto de FL,
realizando experimentos con una CNN de siete capas y ResNet-50 bajo condiciones IID para identificar
qué capas exhibían similitud inter-cliente natural.
Sus resultados confirmaron que las capas tempranas convergían espontáneamente a representaciones
similares entre clientes incluso sin regularización explícita, mientras que las capas finales
---especialmente el clasificador--- mostraban alta disimilitud inter-cliente.
A partir de esta observación, propusieron FedCKA, un método que añade una penalización de
regularización basada en CKA únicamente sobre las capas naturalmente similares, con una eficiencia
computacional casi constante respecto de la profundidad de la red.
El trabajo estableció dos contribuciones metodológicas directamente relevantes para esta tesis:
la validación de CKA como métrica de divergencia representacional en FL y la constatación empírica
de que el drift no es uniforme a lo largo de la profundidad de la red.
Sin embargo, el análisis se realizó solo bajo IID como línea base y únicamente sobre arquitecturas
CNN, sin explorar cómo varían estos patrones entre familias arquitectónicas distintas ni bajo
diferentes niveles de heterogeneidad no-IID.

La base teórica del CKA lineal fue formalizada por Kornblith et al.\ \cite{DBLP:journals/corr/abs-1905-00414},
quienes demostraron que la métrica CKA es invariante a transformaciones ortogonales e isotropías de escala, propiedad fundamental para comparar representaciones aprendidas de forma independiente por distintos clientes.
Esta invarianza hace del CKA lineal el candidato natural para comparar las representaciones
internas de los modelos locales entre sí y respecto del modelo global a lo largo de las rondas de
comunicación, sin que las diferencias métricas sean artefactos de distintas bases de representación.

Desde la perspectiva del análisis de divergencia paramétrica y de gradientes, Li et al.\
introdujeron en el contexto de FedPVR dos diagnósticos de gran relevancia: la
\textit{drift diversity}
$\xi^r = \sum_i \|\mathbf{m}_i^r\|^2 / \|\sum_i \mathbf{m}_i^r\|^2$
---que mide el grado en que las actualizaciones de los clientes apuntan en direcciones
contradictorias--- y la similitud CKA por capa a lo largo del entrenamiento \cite{son2022comparisons}.
Empleando ambas métricas sobre VGG-11 con CIFAR-10 bajo dos niveles de heterogeneidad Dirichlet
($\alpha = 100$ e $\alpha = 0.1$), los autores demostraron que las capas de extracción de
características mostraban invarianza a la heterogeneidad, mientras que las capas profundas
---especialmente el clasificador--- exhibían fuerte sensibilidad.
Esta asimetría por profundidad sustentó el algoritmo FedPVR, que aplica corrección de varianza
solo sobre el clasificador, obteniendo convergencia más rápida que FedAvg, FedProx y SCAFFOLD
con menor sobrecarga de comunicación.
El trabajo validó metodológicamente la combinación de CKA con similitud de gradientes por capa
como herramienta de diagnóstico, aunque su alcance quedó restringido a una sola arquitectura CNN
y no consideró la variación de la asimetría entre familias arquitectónicas distintas.

Yasar et al.\ propusieron L-DAWA, un método de agregación por capas que ponderaba la contribución
de cada capa de cada cliente mediante la divergencia angular (similitud coseno) entre los
parámetros locales y el modelo global: $\delta_k^{(l)} = (\mathbf{w}_g^{(l)} \cdot \mathbf{w}_k^{(l)}) /
(\|\mathbf{w}_g^{(l)}\| \|\mathbf{w}_k^{(l)}\|)$ \cite{son2022comparisons}.
Sus experimentos sobre ResNet-18 en aprendizaje auto-supervisado mostraron que la divergencia
angular variaba drásticamente entre capas y no podía ser capturada por un escalar a nivel de
modelo, lo que apoya el diseño de esta investigación de registrar métricas por capa de forma
continua a lo largo de las rondas.

Los trabajos de personalización por capas en FL aportaron una perspectiva adicional.
Tan et al.\ demostraron mediante un hipernetwork por cliente (pFedLA) que los pesos de agregación
óptimos eran distintos para cada capa, confirmando que la utilidad de cada capa para el modelo
global variaba de forma heterogénea \cite{son2022comparisons}.
Nguyen et al.\ mostraron con FedLAG que los conflictos de gradiente entre capas no seguían un
patrón monótono de entrada a salida ---como asumían trabajos anteriores--- sino una distribución
compleja determinada por el tipo de capa y el grado de heterogeneidad, y que la personalización
selectiva de las capas con mayor conflicto producía mejoras de 5--40\% sobre los métodos de
agregación uniforme bajo alta heterogeneidad \cite{DBLP:journals/corr/abs-1910-06378}.

En conjunto, la literatura sobre análisis por capas ha demostrado que: (1) el drift es
estructuralmente no uniforme a lo largo de la profundidad de la red; (2) las capas de normalización
y clasificación son sistemáticamente las más susceptibles bajo no-IID en CNNs; y (3) CKA y la
similitud coseno de gradientes son instrumentos complementarios para caracterizar la divergencia
representacional y paramétrica, respectivamente.
No obstante, todos los estudios revisados se circunscribieron a arquitecturas CNN.
La extensión de estos diagnósticos a ViTs, SSMs y Vision GNNs ---donde la noción de capa de
características versus clasificador no es directamente equiparable--- permanece como un vacío no
explorado, que constituye la contribución metodológica central de esta tesis.


\section{Normalización en FL heterogéneo}

La relación entre la estrategia de normalización y el comportamiento del modelo bajo heterogeneidad
estadística fue formalizada teóricamente por Wang et al.\ \cite{wang2023batch}, quienes demostraron
analíticamente que BatchNorm introduce un sesgo de gradiente en FL no-IID: cada cliente computa
estadísticas de media y varianza sobre un mini-batch sesgado respecto a la distribución global,
lo que produce correcciones de normalización inconsistentes que se propagan durante la agregación
y corrompen el modelo global.
Du et al.\ \cite{du2022rethinking} extendieron este análisis de forma empírica, comparando BN,
LN, GN e IN sobre múltiples configuraciones federadas y concluyendo que las estrategias basadas
en la muestra individual ---LN y GN--- ofrecían mayor estabilidad.
Casella et al.\ \cite{casella2023experimentingnormalizationlayersfederated} confirmaron este
resultado de forma sistemática: LN y GN mantuvieron precisión competitiva bajo particiones no-IID
donde BN degradaba sensiblemente el rendimiento, mientras que el impacto del número de grupos en
GN era secundario frente a la dicotomía entre normalización por batch y por muestra.

La formulación original de Group Normalization \cite{DBLP:journals/corr/abs-1803-08494} propuso
dividir los canales de activación en grupos y normalizar dentro de cada grupo, eliminando la
dependencia estadística entre muestras del batch.
Esta propiedad hace de GN una alternativa directamente sustituible en redes convolucionales:
conserva el sesgo inductivo de la convolución mientras elimina la principal fuente de inestabilidad
identificada bajo FL heterogéneo.
GN con un único grupo equivale a Layer Normalization, caso límite de máxima independencia
respecto al batch.

La síntesis de esta línea muestra consenso en torno a la superioridad de LN y GN sobre BN en
condiciones no-IID.
Sin embargo, ninguno de los estudios citados ha medido cómo el reemplazo de BatchNorm modifica
específicamente el perfil de divergencia de parámetros por capa a lo largo de las rondas de
comunicación, ni ha comparado ese perfil modificado con el de arquitecturas que emplean LN de
forma nativa (ViT, Vim, ViG).
Esta ablación controlada ---que la presente investigación realiza sobre EfficientNet-B1---
constituye la pieza metodológica que permite aislar la contribución de la normalización respecto
del sesgo inductivo convolucional como fuentes independientes de inestabilidad federada.


\section{Síntesis y posicionamiento de la investigación}

La revisión de la literatura permite trazar un mapa preciso de lo que se conoce y lo que permanece
abierto.
La arquitectura es un moderador primario del \textit{client drift} bajo no-IID
\cite{qu2022rethinking, pieri2023handlingdataheterogeneityarchitectural, buyuktacs2024transformer}:
las familias basadas en normalización por muestra y mecanismos de integración global son
sistemáticamente más robustas que las CNNs con BatchNorm, y esta diferencia supera en magnitud a
la producida por la elección del algoritmo de agregación.
BatchNorm ha sido identificada como el principal vector de inestabilidad federada en CNNs
\cite{wang2023batch, du2022rethinking, xu2023fedconv, siomos2024architecture}, aunque la evidencia
disponible no desagrega su contribución respecto del sesgo inductivo convolucional en términos de
divergencia de parámetros por capa.

El análisis por capas ha demostrado que el drift tiene estructura interna compleja: se concentra
en las capas de clasificación y normalización, mientras que las capas de extracción de bajo nivel
muestran mayor estabilidad relativa \cite{son2022comparisons}.
CKA y la similitud coseno de gradientes inter-cliente han emergido como las métricas más
informativas para caracterizar esta estructura, complementando las medidas paramétricas con
información sobre la divergencia representacional funcional
\cite{DBLP:journals/corr/abs-1905-00414, son2022comparisons}.
No obstante, todos estos diagnósticos han sido validados exclusivamente sobre CNNs.

Las familias SSM han demostrado viabilidad técnica en contextos federados \cite{wu2024mamba} y
competitividad en dominios médicos \cite{yue2024medmamba}, pero ningún estudio ha caracterizado
su drift bajo condiciones no-IID controladas.
Vision GNN ha sido integrado en FL por primera vez en FedVIG \cite{cui2025fedvig}, pero sin
análisis de heterogeneidad sistemática.
No existe ningún estudio que haya comparado simultáneamente las cuatro familias ---CNN, ViT, SSM,
Vision GNN--- bajo un protocolo experimental controlado con métricas de drift por capa.

La presente investigación cierra estas brechas mediante: (1) una evaluación transversal de las
cuatro familias arquitectónicas bajo particiones de Dirichlet con cinco niveles de heterogeneidad
sobre Brain Tumor MRI y CIFAR-10; (2) el registro continuo de divergencia de pesos ($L_2$),
alineación de gradientes (similitud coseno) y divergencia representacional (CKA) por tipo de capa
y por ronda de comunicación; y (3) una ablación de normalización sobre EfficientNet-B1 que disocia
la contribución de BatchNorm del sesgo inductivo convolucional como fuentes independientes de
inestabilidad federada.
Este diseño es el primero en combinar la amplitud de familias arquitectónicas con la profundidad
del análisis por capa y el rigor de una parametrización continua de la heterogeneidad,
posicionando la investigación como una extensión natural y necesaria del estado del arte revisado.
